\section{Introduction}
\label{sec:intro}

Language models increasingly act rather than answer. They choose moves in games, calls to tools and steps in workflows, and what they produce is executed instead of read \citep{yao2023react,liu2024agentbench,paglieri2025balrog}. Many of these systems run on open weights, and the small end of the open-weight range is being promoted for agentic work because it is cheaper, faster and can run locally \citep{belcak2025slm}. Anyone building such a system faces a practical question: how small can the model be before its decisions get worse, and what can be done about it when they do?

The literature gives mixed signals. Pre-training loss falls smoothly with parameter count \citep{kaplan2020scaling,hoffmann2022chinchilla}, but decision quality need not. Some abilities appear abruptly at scale \citep{wei2022emergent}, although many such jumps turn out to be artifacts of discontinuous metrics \citep{schaeffer2023mirage}. Several failures persist as models grow. Language models explore poorly in multi-armed bandits unless the history is summarized for them \citep{krishnamurthy2024explore}. Models from 2B to 27B parameters leave much of the action space untried and often act against their own correct rationale \citep{schmied2026greedy}, and they can fail even to exploit what they have already seen \citep{harris2026explore}. Other results favor small models. A 3B model post-trained with multi-step GRPO beat a 72B baseline on Frozen Lake \citep{dilkes2025reinforced}, and distilling bandit algorithms into smaller models let them explore better than larger models without that help \citep{nie2025evolve}. Most agent benchmarks evaluate frontier APIs or heterogeneous collections of open models \citep{wu2024smartplay,liu2024agentbench,paglieri2025balrog,ruoss2025lmact,duan2024gtbench}, so when one model beats another they cannot tell whether size or training recipe made the difference.

There is a second, less discussed variable. A language model is a distribution over text, not a policy, and something has to turn one into the other. The usual choice is to generate text and parse an action out of it, which mixes formatting failures with bad decisions \citep{huang2022zeroshot,liu2024agentbench}. One can instead score every legal action by its likelihood \citep{carta2023glam,tan2024twosome}, read the decision off the hidden states with a probe \citep{alain2016probes,orgad2025know}, or fine-tune on expert actions \citep{chu2025sft}. These methods use the weights in different ways, and in multiple-choice question answering they disagree with each other more often than one would hope \citep{robinson2023mcsb,wang2024myanswerc,wang2024lookattext}. Papers on them usually study one method at one or two model sizes. A small model that fails when prompted might succeed when scored or probed; we do not know how often that happens, or whether it depends on the kind of decision.

The closest existing studies each cover part of this space. \citet{jiwatode2026spatial} vary Qwen3 model size, thinking mode and planning horizon on three spatial-navigation games, with one model family and prompting as the only conversion. \citet{schmied2026greedy} study models from 2B to 27B on bandits, contextual bandits and tic-tac-toe, identify greediness, frequency bias and a knowing-doing gap, and reduce all three with RL fine-tuning on the model's own reasoning; their comparison is between prompting and RL fine-tuning, not across a set of conversion methods on a size ladder. \citet{monea2025icrl} span 500M to 70B parameters, but for in-context bandit learning on classification tasks, and \citet{nie2025evolve} relate bandit regret to model size and test inference-time and distillation interventions. None of these crosses same-recipe model ladders with several conversion methods and several kinds of decision problem, normalizes against an oracle, or fixes its hypotheses before collecting data. Some of our hypotheses already have partial answers in this work, and Section~\ref{sec:related} says where. For those, what we add is a controlled, within-family and pre-registered measurement, not a new direction of effect.

Our contributions are the following.
\begin{itemize}
  \item \textbf{A benchmark.} Six procedurally generated decision tasks, each with an exact or canonical oracle, covering exploration (a multi-armed bandit), in-context rule learning (a contextual bandit), one-shot tabular decisions under covariate shift and under a reversed effect (loan approval), deterministic planning (gridworld), adversarial play (tic-tac-toe) and risk (blackjack). A seed fixes an instance and all of its randomness, so policies are compared on identical potential outcomes, and returns are mapped onto a scale on which the random policy scores 0 and the oracle 1 (Section~\ref{sec:benchmark}).
  \item \textbf{A harness.} Five conversion methods behind one interface: free-form generation, exact likelihood scoring of legal actions, a logistic probe on hidden states, LoRA behavior cloning, and the same probe applied to raw task features as a control that uses no language model. The harness runs local Hugging Face models and models behind any OpenAI-compatible server, logs every decision with the raw model output, and regenerates every result table and figure in this paper from those logs.
  \item \textbf{A pre-registered study.} Within-family ladders (Qwen2.5 from 0.5B to 72B, Qwen3, SmolLM2 and Gemma 3, with Pythia and OLMo 2 as fully open controls), paired seeds, and seven hypotheses whose falsification criteria are frozen at a git tag before the main sweep (Section~\ref{sec:setup}).
  \item \textbf{Findings.} \result{[One sentence per hypothesis, for example: the slope under generation is $b$ per tenfold increase in parameters and is steeper on planning tasks (H1); a LoRA-tuned model of at most 1.7B matches a ten times larger zero-shot model on $k$ of six tasks in distribution but not under shift (H3); scoring closes $Z$\% of the small-to-large gap (H4); 4-bit weights cost $q$ below 2B (H5); no prompted model reaches UCB1 on the bandit (H6); chain-of-thought helps planning only from $m$B upward (H7)]} (Sections~\ref{sec:results} and~\ref{sec:analysis}).
\end{itemize}

Section~\ref{sec:related} reviews related work and marks where it already answers parts of our hypotheses. Section~\ref{sec:benchmark} defines the tasks, the score and the conversion methods, and Section~\ref{sec:setup} the models, factors and statistics. Section~\ref{sec:results} tests each hypothesis against its pre-registered criterion. Section~\ref{sec:analysis} examines errors, cost and the language-model-free control, and Section~\ref{sec:limitations} sets out what the design cannot tell us. The appendices give every prompt verbatim, the task generators, the full pre-registration, the compute budget, a reproducibility checklist and the verified model identifiers.
